% Book Two Pass 2: planar element derivation and sign/measure audit.
\clearpage
\section*{Research reconstruction: from triangle gradients to a Laplacian}
\noindent\textbf{Primary lineage:} the finite-element/cotangent discrete
Laplacian described in geometric processing literature, the
\href{https://libigl.github.io/tutorial/}{libigl tutorial}, and
Nicholas Sharp and Keenan Crane,
\href{https://www.cs.cmu.edu/~kmcrane/Projects/NonmanifoldLaplace/index.html}
{\emph{A Laplacian for Nonmanifold Triangle Meshes} (2020)}.

\paragraph{Where the matrix comes from.}
Consider a triangle of area \(A>0\) with piecewise-linear hat
basis functions \(\phi_0,\phi_1,\phi_2\). Each \(\phi_i\) equals one
at vertex \(i\) and zero at the other two vertices. Its gradient is
constant on the triangle, because each basis function is linear.
For vertex values \(u_i\), the interpolated field and gradient are
\[
u(x)=\sum_i u_i\phi_i(x), \qquad
\nabla u=\sum_i u_i\nabla\phi_i.
\]
The integrated squared-gradient energy is
\[
E(u)=\frac12\int_T\lVert\nabla u\rVert^2\,dA
=\frac12u^TK_Tu,\qquad
(K_T)_{ij}=A\,\nabla\phi_i\cdot\nabla\phi_j.
\]
Assembly over all triangles yields a symmetric
positive-semidefinite stiffness \(K\), with \(K\mathbf 1=0\)
on a mesh without fixed-value boundary constraints: a constant
field has zero gradient and therefore zero energy. A lumped
vertex mass for this simple P1 experiment is
\(M_{ii}=\sum_{T\ni i}A_T/3\).
This is an \emph{area-weighted measure}, not the same object as
the stiffness matrix.

\paragraph{An exact numerical triangle.}
Take \(v_0=(0,0)\), \(v_1=(2,0)\), \(v_2=(0.5,0.5)\).
The area is \(A=0.5\), and the elemental matrix is
\[
K_T=
\begin{pmatrix}
1.25&0.25&-1.5\\
0.25&0.25&-0.5\\
-1.5&-0.5&2
\end{pmatrix}.
\]
It is symmetric, has zero row sums, and satisfies
\(u^TK_Tu\ge0\) for every \(u\).
Yet \(K_{01}=+0.25\) is a positive off-diagonal entry. An
energy-stable discrete operator is not therefore guaranteed to
be a nonnegative-weight averaging operator or to satisfy a
maximum principle. Mesh quality, angle and weighting conventions
are material to the result.

\paragraph{Why nonmanifold geometry is a separate paper.}
The standard cotangent stencil assumes a suitable triangulated
domain and imposes no automatic guarantee that every input
mesh has the desirable spectral, positivity or boundary properties.
The Sharp--Crane paper develops a \emph{tufted cover} to extend
Laplacian construction to more general triangle meshes. A basic
P1 matrix on an ordinary planar grid does \emph{not} reproduce that
construction, its intrinsic triangulation strategy or its results.

\paragraph{Reconstruction.}
Use \texttt{tools/labs/cotan\_energy\_lab.py} for the single-element
witness and \texttt{tools/labs/book02\_papers\_numeric.py} for an
assembled triangular-grid stiffness/mass system. Check zero
row sum, area conservation, energy positivity, exact linear-field
Dirichlet energy, and degenerate triangle rejection.
After this baseline is correct, study boundaries and nonmanifold
cases as a separate method rather than assuming the original
cotangent operator can simply be reused unchanged.

\paragraph{Sign warning.}
Some software exposes a Laplacian \(L\) with negative diagonal
entries, while this laboratory uses the positive-semidefinite
stiffness \(K=-L\). Always write the evolution equation and
solver's sign convention before comparing matrix entries.
